% ===========================================================================
% Περίληψη - ΓΡΑΨΤΗΝ ΤΕΛΕΥΤΑΙΑ, μαζί με το \cref{ch:greek-summary}.
%
% Μία σελίδα. Η δομή που δουλεύει: πρόβλημα -> προσέγγιση -> τι χτίστηκε ->
% ΤΙ ΜΕΤΡΗΘΗΚΕ (με αριθμούς) -> τι δεν έγινε.
% Η τελευταία πρόταση να είναι ειλικρινής για το όριο της εργασίας.
%
% Προσοχή: κάθε λατινική λέξη μέσα σε ελληνικό κείμενο θέλει \en{...}.
% ===========================================================================
\chapter*{Περίληψη}
\addcontentsline{toc}{chapter}{Περίληψη}
\selectlanguage{greek}

% TODO

\vspace{1em}
% \elkeywords is provided by the template.
\elkeywords{}
% π.χ.: συνελικτικά νευρωνικά δίκτυα, \en{MobileNetV2}, κβαντοποίηση
% \en{int8}, \en{FPGA}, \en{Zynq UltraScale+},
% επιταχυντής υλικού, \en{Verilog}, επαλήθευση ακριβείας bit,
% ανίχνευση ασθενειών αμπέλου

\selectlanguage{english}
\cleardoublepage
